\documentclass{article}

\usepackage[utf8]{inputenc}
\usepackage[T1]{fontenc}
\usepackage[spanish]{babel}
\usepackage{amsmath}
\usepackage{amssymb}
\usepackage{booktabs}
\usepackage{geometry}
\geometry{margin=2.5cm}

\title{Modelado estadístico de las órdenes médicas}
\date{}

\begin{document}

\maketitle

\section{Grado de criticidad}

La frecuencia con la que se generan nuevas órdenes médicas y el caudal
indicado en ellas dependen del estado clínico del paciente.
Pacientes que presentan un cuadro crítico requieren revisiones frecuentes y
ajustes continuos de la infusión, mientras que pacientes estables suelen
mantener la misma configuración durante períodos más prolongados.

Con el objetivo de representar este comportamiento, se define el parámetro
$k$, denominado \textbf{grado de criticidad del paciente}:
\[
  k \in [0,1]
\]
donde $k=1$ corresponde a un paciente crítico y $k=0$ a un paciente estable.
Este parámetro influye tanto en el tiempo transcurrido entre órdenes médicas
como en el caudal indicado por cada orden.

% ─────────────────────────────────────────────────────────────
\section{Tiempo entre órdenes médicas}

Se define la variable aleatoria discreta $T$ como el tiempo transcurrido,
expresado en horas, entre dos órdenes médicas consecutivas:
\[
  T \in \{2,\,4,\,6,\,8,\,12\}
\]

\subsection*{Distribuciones extremas}

\begin{table}[h]
\centering
\begin{tabular}{@{}ccc@{}}
\toprule
$t$ (horas) & $P_C(T=t)$ & $P_E(T=t)$ \\
\midrule
2  & 0.50 & 0.00 \\
4  & 0.30 & 0.05 \\
6  & 0.15 & 0.15 \\
8  & 0.05 & 0.40 \\
12 & 0.00 & 0.40 \\
\bottomrule
\end{tabular}
\caption{Distribución crítica $P_C$ y distribución estable $P_E$ del tiempo entre órdenes.}
\end{table}

\subsection*{Distribución general}

La distribución paramétrica en función de $k$ se define como:
\[
  P(T=t \mid k) \;=\; k\,P_C(T=t) + (1-k)\,P_E(T=t).
\]

El tiempo promedio entre órdenes médicas es:
\[
  \mathbb{E}[T \mid k]
  \;=\; \sum_{t \in \{2,4,6,8,12\}} t\;P(T=t \mid k).
\]

% ─────────────────────────────────────────────────────────────
\section{Caudal de infusión}

Se define la variable aleatoria discreta $C$ como el caudal indicado por
una orden médica, expresado en ml/h:
\[
  C \in \{0,\,50,\,100,\,150,\,200\}.
\]
El valor $C=0$ representa la detención de la infusión.

\subsection*{Distribuciones extremas}

En pacientes críticos predominan los caudales elevados; la probabilidad de
emitir una orden con $C=0$ es prácticamente nula.
En pacientes estables los requerimientos son menores y resulta más probable
la suspensión de la infusión.

\begin{table}[h]
\centering
\begin{tabular}{@{}ccc@{}}
\toprule
$c$ (ml/h) & $P_C(C=c)$ & $P_E(C=c)$ \\
\midrule
0   & 0.00 & 0.30 \\
50  & 0.10 & 0.35 \\
100 & 0.20 & 0.20 \\
150 & 0.40 & 0.10 \\
200 & 0.30 & 0.05 \\
\bottomrule
\end{tabular}
\caption{Distribución crítica $P_C$ y distribución estable $P_E$ del caudal.}
\end{table}

\subsection*{Distribución general}

\[
  P(C=c \mid k) \;=\; k\,P_C(C=c) + (1-k)\,P_E(C=c).
\]

En particular, la probabilidad de suspensión resulta:
\[
  P(C=0 \mid k) \;=\; 0.30\,(1-k).
\]
A medida que disminuye la criticidad, aumenta la probabilidad de que una
orden indique la finalización de la infusión.

El caudal medio esperado es:
\[
  \mathbb{E}[C \mid k]
  \;=\; \sum_{c \in \{0,50,100,150,200\}} c\;P(C=c \mid k).
\]

% ─────────────────────────────────────────────────────────────
\section{Funciones generadoras}

A partir de las distribuciones anteriores se definen las funciones
generadoras de muestras:
\[
  \mathrm{Caudal}(k) \sim P(C \mid k),
  \qquad
  \mathrm{Tiempo}(k)  \sim P(T \mid k).
\]
Estas funciones son utilizadas por el generador DEVS de órdenes médicas
(véase el documento de modelos DEVS) para producir, en cada transición
interna, el próximo caudal y el próximo intervalo entre órdenes.

\end{document}